The validation reported in Section~\ref{sec:results} was performed on a laboratory bench prototype.
No clinical deployment took place, no patients participated, and no patient records were collected
or retained.

The system nevertheless handles identifying personal data by design. Records live only in volatile
memory, so the device accumulates no persistent identifiable data, and the storage proposed in
Section~\ref{sec:limitations} reverses that property and would need a retention limit and
encryption at rest. Because the console shows a patient's name beside their vitals in plain text,
its placement on the desk is a confidentiality decision as well as an ergonomic one. The absence of
connectivity is itself a privacy benefit, since no data leaves the device, and any wireless
extension should therefore be treated as a change requiring authenticated pairing, encrypted
transport and explicit patient consent. Finally, the temperature and pulse channels are triage
indicators, accurate enough to draw a physician's attention and not accurate enough to support a
clinical decision alone. Misreading that boundary is the most plausible route from this work to
harm, and the system should not be deployed or extended as an autonomous screening device.
